\begin{theorem}[\texorpdfstring{\hyperref[prf:cantors-intersection-theorem]{Cantor's Intersection Theorem}}{Cantor's Intersection Theorem}]
\label{thm:cantors-intersection-theorem}
Let $(X, d)$ be a metric space and let
\[
  A_1 \supseteq A_2 \supseteq A_3 \supseteq \cdots
\]
be a nested sequence of nonempty compact subsets of $X$. Then
\[
  \bigcap_{n=1}^{\infty} A_n
\]
is compact and nonempty.
\end{theorem}

\begin{remark*}[Fully quantified form]
$\forall (X,d)$ metric space,
$\forall (A_n)_{n \in \mathbb{N}}$ with each $A_n \subseteq X$ compact, $A_n \neq \emptyset$,
and $A_{n+1} \subseteq A_n$:
\[
  \bigcap_{n=1}^{\infty} A_n \neq \emptyset
  \quad\text{and}\quad
  \bigcap_{n=1}^{\infty} A_n \text{ is compact.}
\]
\end{remark*}

\begin{remark*}[Intuition]
Every compact set is closed, so the intersection of any collection of compact sets
is closed. Closedness alone does not prevent the intersection from being empty
(nested open discs can vanish to nothing). Compactness adds the extra force:
it guarantees that a sequence of witnesses — one chosen from each $A_n$ — cannot
escape to infinity or disappear, and must cluster at a point that lands in every set.
\end{remark*}

\begin{remark*}[Consequence]
The sharpened form (Corollary below) shows that if additionally
$\operatorname{diam}(A_n) \to 0$, then the intersection is a single point.
This is the engine behind the Baire Category Theorem and completeness arguments.
\end{remark*}

\begin{remark*}[Common error]
The hypotheses cannot be weakened to closed sets alone. The nested open discs
$A_n = \{\,z \in \mathbb{R}^2 : |z - (1/n, 0)| < 1/n\,\}$
are nonempty, nested, and closed in themselves, but $\bigcap A_n = \emptyset$.
Compactness is essential.
\end{remark*}




\begin{proof}
\Claim Let $(X,d)$ be a metric space and let $A_1 \supseteq A_2 \supseteq \cdots$
be a nested sequence of nonempty compact subsets of $X$. Then
$\bigcap_{n=1}^{\infty} A_n$ is compact and nonempty.

\Given A metric space $(X,d)$ and a nested sequence of nonempty compact sets
$(A_n)_{n \in \mathbb{N}}$ with $A_{n+1} \subseteq A_n$ for all $n$.

\Goal To show that $A := \bigcap_{n=1}^{\infty} A_n$ is compact and nonempty.

\medskip
\noindent\textit{Part 1: $A$ is compact.}

\Strategy We show $A$ is a closed subset of the compact set $A_1$, hence compact.

\ByThm{Every compact set is closed} each $A_n$ is closed in $X$.
\ByThm{Arbitrary intersections of closed sets are closed}
$A = \bigcap_{n=1}^{\infty} A_n$ is closed in $X$.
Since $A \subseteq A_1$ and $A_1$ is compact,
\ByThm{Every closed subset of a compact set is compact}
$A$ is compact. \checkbox

\medskip
\noindent\textit{Part 2: $A$ is nonempty.}

\Strategy We construct a sequence with one point drawn from each $A_n$, extract
a convergent subsequence using compactness of $A_1$, and show its limit lies in
every $A_n$.

Since each $A_n$ is nonempty, \Let $(a_n)_{n \in \mathbb{N}}$ be a sequence with
$a_n \in A_n$ for every $n$.

Because the sets are nested, $a_n \in A_n \subseteq A_1$ for all $n$,
so the entire sequence $(a_n)$ lies in $A_1$.
\ByDef{sequential compactness of $A_1$} there exists a subsequence
$(a_{n_k})_{k \in \mathbb{N}}$ and a point $p \in A_1$ such that
\[
  a_{n_k} \to p \quad \text{in } (X,d).
\]

\Let $m \in \mathbb{N}$ be arbitrary. For all $k$ sufficiently large, $n_k \geq m$,
so $a_{n_k} \in A_{n_k} \subseteq A_m$. \Therefore the tail of the subsequence
$(a_{n_k})$ lies in $A_m$.
\ByThm{Closed sets contain the limits of their convergent sequences}
and since $A_m$ is compact (hence closed), we conclude $p \in A_m$.

Since $m$ was arbitrary, $p \in A_m$ for every $m \in \mathbb{N}$,
hence $p \in \bigcap_{n=1}^{\infty} A_n = A$. \checkbox

\Hence $A$ is nonempty. \AsReq
\end{proof}

\begin{remark*}[Proof shape]
Two-part direct proof. Part 1 is short: the intersection of closed sets is closed,
and a closed subset of a compact set is compact. Part 2 is the heart: pick one
point from each $A_n$, forming a sequence in $A_1$; extract a convergent
subsequence via compactness of $A_1$; argue that the limit lands in every $A_m$
by the tail argument and closedness of each $A_m$. The nesting is used exactly
at the step $A_{n_k} \subseteq A_m$ for large $k$.
\end{remark*}

\begin{remark*}[Dependencies]
\begin{itemize}
  \item Definition of sequential compactness: \ref{def:sequential-compactness}
  \item Every compact set is closed: \ref{prop:compact-implies-closed}
  \item Arbitrary intersections of closed sets are closed: \ref{thm:intersection-closed-sets-closed}
  \item Every closed subset of a compact set is compact: \ref{prop:closed-subset-compact}
  \item Closed sets contain limits of convergent sequences (sequential characterization): \ref{def:closed-set-sequential}
\end{itemize}
\end{remark*}























